\documentclass[10pt]{article}
\usepackage[margin=20mm]{geometry}
\usepackage{amsmath,amssymb,mathtools,fontspec,unicode-math,xurl,hyperref}
\IfFontExistsTF{DejaVu Serif}{\setmainfont{DejaVu Serif}}{\setmainfont{Latin Modern Roman}}
\setmathfont{latinmodern-math.otf}
\setlength{\parindent}{0pt}\setlength{\parskip}{6pt}\sloppy
\title{Technical verification of an almost sharp quantitative Duffin–Schaeffer theorem}
\author{AI-assisted technical review}
\date{}
\begin{document}
\maketitle
Abstract: A claim-focused examination of the variance reduction, overlap estimates and GCD-graph argument identifies one local defect: the strict quality-increment alternative in Lemma 10.1 fails when the permitted measure has zero-weight vertices. An explicit counterexample and a positive-support repair are given. The repair preserves the graph quantities and restores the high-degree subgraph lemma used downstream. No failure of Theorems 1 or 2 is established. The separate passage from a variance asymptotic to almost-everywhere sharpness requires an additional argument if Section 16 is intended as a standalone proof. Scientific adjudication remains pending.\par
\section*{Identity, reviewed version and scope}
This report concerns Dimitris Koukoulopoulos, James Maynard and Daodao Yang, An almost sharp quantitative version of the Duffin–Schaeffer conjecture, assigned identifier M25-0161. The exact reviewed source is arXiv:2404.14628v2, bearing the arXiv date 20 September 2024 and an internal date of 23 September 2024 [source]. Its 41 PDF pages have printed page numbers 1–41; all page references below are one-based PDF pages, equal to the printed numbers. The assignment links the article to DOI 10.1215/00127094-2024-0055 (https://doi.org/10.1215/00127094-2024-0055). The acquired author manuscript is not asserted to be the journal version of record.\par
The examination covered the statements and dependency chain in Sections 1–16, with independent algebraic reconstruction of the dyadic variance argument, the anatomical estimates, the main graph-quality inequalities and the prime-supported variance calculation. Three native in-chat assistants checked disjoint portions of Sections 8–16; the lead reviewer independently checked the reported defect and its repair. The retained coverage ledger distinguishes direct checks, helper checks and external results not re-proved. In particular, cited sieve and multiplicative-function theorems and the historical central limit theorem were not independently verified from their original publications. No separate supplement or released repository was available in the assigned source package. No author code was run.\par
\section*{Claims and checks}
Let \(0\leq\psi(q)\leq 1/2\), and set \[\Psi(Q)=2\sum_{q\leq Q}\frac{\varphi(q)\psi(q)}q,\qquad
 N(\alpha;Q)=\#\{(a,q)\in\mathbb Z\times\mathbb N:
 q\leq Q,\ (a,q)=1,\ |\alpha-a/q|<\psi(q)/q\}.\] Theorem 1 assumes \(\Psi(Q)\to\infty\) and gives the almost-everywhere error \(O_\varepsilon(\Psi(Q)^{1/2+\varepsilon})\). Theorem 2 bounds the second centered moment by \(\Psi(Q)+O_\varepsilon(\Psi(Q)^{1+\varepsilon})\). The review did not find a counterexample to either statement.\par
\subsection*{Normalization and the variance-to-pointwise reduction}
The disjoint intervals on page 3 give \(\lambda(A_q)=2\varphi(q)\psi(q)/q\); the boundary denominator \(q=1\) has two truncated intervals whose combined length is \(2\psi(1)\). Thus \(N=\sum_{q\le Q}1_{A_q}\) and \(\mathbb EN=\Psi\). Since each increment of \(\Psi\) is at most one, the integer cutoff \(Q_j\) in Section 4 satisfies \(j-1\leq\Psi(Q_j)<j\). For a binary block of length \(2^s\), Theorem 2 applied to the truncated function \(1_{q>Q'}\psi(q)\) gives a second-moment bound \(O_\varepsilon(2^{s(1+\varepsilon)})\). There are at most \(2^{r-s}\) such blocks at scale \(s\). Consequently the exceptional-set estimate is \[\lambda(E_r)\ll_\varepsilon
 (r+1)2^{-(1+2\varepsilon)r}
 \sum_{s=0}^r 2^{r-s}2^{s(1+\varepsilon)}
 \ll_\varepsilon(r+1)2^{-\varepsilon r}.\] This is summable. A countable sequence of positive \(\varepsilon\)’s yields a common full-measure set for every positive exponent allowance. The factor-of-two normalization in the first display on page 7 does not affect this bound: the actual block expectation is \(2\sum\psi(q)\varphi(q)/q=\Psi(Q)-\Psi(Q')\).\par
\subsection*{Overlap, anatomical estimates and dyadic sums}
For Lemma 5.1, the triangular weight has integral \(2\delta\Delta\) and total variation \(2\delta\). Hence the Riemann-sum error in (5.3) is uniformly \(O(\delta)\). The coprimality relaxation to the \(t\)-smooth part, the square-free divisor expansion and the estimate \(\prod_{p>t}(1-1/p)^{-1}\leq
\exp(2\sum_{p>t}1/p)\) give the stated overlap upper bound, conditional on the initial exact overlap-counting identity.\par
The exponential moments in Lemmas 6.2–6.3 are compatible with Lemma 6.1: one chooses a fixed divisor-function order large enough to dominate the prime-power values. The large-prime contribution is bounded using \(t\sum_{p>t}p^{-2}=O(1)\); the small-prime moment grows at most as a fixed power of \(\log(2t)\), absorbed into the spare factor \(t\). These bounds explain the two exceptional-pair savings.\par
The Section 7 sums have the required decay: \(2^{i(\varepsilon/2-1)}2^{i(1-\varepsilon)}
=2^{-i\varepsilon/2}\), and the anatomical exceptional classes have summable exponential or geometric tails. At \(\Psi(Q)<1\), the large-\(D\) remainder may be bounded directly by \(\Psi(Q)^2\); this closes the small-mass case without reversing the order of powers of a number below one. When \(1/2\leq D<1\), replacing the sieve cutoff \(D^2\) by 1 leaves the relevant prime sums unchanged.\par
\subsection*{Graph reductions}
The graph checks followed the modified exponent \(2+\tau\), with \(0<\tau<1/100\), through the quality ratios, density restrictions, anatomical pruning and small-/large-prime iteration. The finite sets of unaccounted gcd primes provide a termination measure. Several compressed arguments can be reconstructed with the changed exponent rather than assumed from the older graph argument. The local zero-weight defect below must be addressed when using maximality to deduce pointwise degree bounds. Exact prime valuations in Proposition 8.1(e) and Lemma 13.2(v) are interpreted on the output vertex sets; the preceding discussion and construction support that reading. No separate substantive finding is assigned to that notation.\par
\section*{A local defect and its verified repair}
Problem 1 (Zero-weight vertices invalidate Lemma 10.1 as stated).  Definition 1(a), page 13, permits a measure on \(\mathbb N\) with finite singleton masses, including zero. Lemma 10.1, pages 22–23, asserts that every GCD graph of positive edge density either has the stated positive weighted degree at every vertex or admits a subgraph with the same multiplicative data and strictly larger quality. These alternatives need not exhaust the allowed graphs.\par
\subsection*{Complete counterexample}
Take \(\tau=1/200\), \(\mu(\{1\})=1\), and \(\mu(\{n\})=0\) for \(n\ne1\). Let \[V=\{1,2\},\qquad W=\{1\},\qquad
 E=\{(1,1)\},\qquad P=\varnothing,\] with both multiplicity functions empty. The vertex sets are finite positive integers; \(E\subseteq V\times W\); all singleton masses are finite; and the prime constraints are vacuous. Furthermore \[\mu(V)=\mu(W)=\mu(E)=1,\qquad \delta(G)=1,\qquad q(G)=1.\] Thus every hypothesis of the lemma holds. The neighborhood of \(2\in V\) is empty, so \[\mu(\Gamma_G(2))=0<
 \frac{1+\tau}{2+\tau}\delta(G)\mu(W)
 =\frac{201}{401}.\] Alternative (a) fails.\par
Any subgraph with the same multiplicative data either retains the sole positive-weight edge \((1,1)\), in which case its density and quality are both 1, or has zero edge measure and quality 0. Therefore no such subgraph has quality larger than 1, and alternative (b) also fails. This is an exact finite argument, not a floating-point test. The independent finite enumeration retained with the helper evidence agrees with it.\par
\subsection*{Where the proof fails}
On page 23 the proof multiplies the strict low-degree inequality by \(\mu(v)\) and continues with a strict inequality. That operation preserves strictness only for \(\mu(v)>0\). For the counterexample vertex \(v=2\), deleting it leaves both edge measure and quality unchanged. This is not an extraction artifact: the definition, lemma and strict inequalities were checked on rendered source pages.\par
\subsection*{Repair and exact downstream scope}
Restrict the initial graph to \[V_+=\{v\in V:\mu(v)>0\},\qquad
 W_+=\{w\in W:\mu(w)>0\},\qquad
 E_+=E\cap(V_+\times W_+).\] All removed edges have product measure zero. Consequently \(\mu(V_+)=\mu(V)\), \(\mu(W_+)=\mu(W)\), \(\mu(E_+)=\mu(E)\), and both density and quality are unchanged. The GCD constraints and multiplicative data are inherited. Now choose a maximal subgraph of this positive support. The strict-increment proof applies to every remaining vertex and yields the degree conclusion of Lemma 8.6, with quality at least the original quality, exactly as that lemma requires.\par
For completeness, write \(x=\mu(v)/\mu(V)\in(0,1)\) and \(c=(1+\tau)/(2+\tau)\). Removing a vertex whose degree is below the threshold gives \[\frac{q(G')}{q(G)}>
 \frac{(1-cx)^{2+\tau}}{(1-x)^{1+\tau}}>1,\] because the logarithm on the right vanishes at \(x=0\) and has derivative \[(1+\tau)\left(\frac1{1-x}-\frac1{1-cx}\right)>0.\] This verifies the repaired strict-increment argument. An equivalent formulation allows a proper vertex deletion with nondecreasing quality, using decreasing vertex count for termination.\par
\subsection*{Impact}
The supported impact is Minor: the printed auxiliary dichotomy is false on an allowed degenerate input. The retained positive-support repair restores the downstream high-degree subgraph existence result without changing any weighted quantity. The counterexample does not invalidate the bilinear estimates, the variance theorem or the almost-everywhere counting theorem. The classification reflects that bounded consequence, rather than the length of the repair.\par
\section*{Scope of the separate sharpness argument}
Section 16, pages 38–39, correctly reduces the prime-supported choice \(\psi(q)=1_{q\ {\rm prime}}/q\) to pair-correlation errors \(O(\max(q,r)^{-2})\). Their sum over distinct primes is bounded, as is \(\sum_q\lambda(A_q)^2\). This gives the stated \(\operatorname{Var}N(\alpha;Q)=\Psi(Q)+O(1)\).\par
The introductory assertion that this alone proves almost-everywhere optimality needs an additional passage from second moments to sample behavior. Divergent second moments can be concentrated on shrinking exceptional sets. This remains true for nondecreasing indicator counts: Appendix A gives an exact countermodel. It is a countermodel to that general inference, not a Duffin–Schaeffer counterexample.\par
The source separately invokes a central limit theorem on page 2. That is relevant counterevidence against any claim that the paper’s sharpness conclusion is false. The historical theorem was not independently checked here. Accordingly this review records the standalone implication in Section 16 as Needs Evidence, and does not raise the supported severity of the paper or assert failure of its main theorem.\par
\section*{Conclusion and recorded states}
Dimension | Recorded state\par
Supported issue | Minor: Lemma 10.1 with zero-weight vertices\par
Repair verification | Analytic positive-support repair; Lemma 8.6 restored\par
Open scope question | Needs Evidence: standalone \(L^2\)-to-a.e. sharpness implication\par
Main conclusions | No counterexample established in the checked scope\par
Scientific adjudication | Pending\par
Human review | No human review supplied\par
Author adoption | Unknown; no adoption or author contact requested\par
Later author versions, journal-final changes and errata were not compared. The finding is attached only to the exact version identified above; absence of a later correction is not claimed. No new experiment is required to check the local counterexample or its repair.\par
1 D. Koukoulopoulos, J. Maynard and D. Yang, An almost sharp quantitative version of the Duffin–Schaeffer conjecture, arXiv:2404.14628v2 (2024), https://arxiv.org/abs/2404.14628v2. Exact acquired author manuscript; 41 pages.\par
\section*{A countermodel to the second-moment implication}
Set \(T_1=a_1=1/2\). For \(n\geq2\), define \[x_n=n-a_{n-1},\quad
 T_n=x_n-\sqrt{x_n^2-1}
     =\frac1{x_n+\sqrt{x_n^2-1}},\quad
 a_n=a_{n-1}+T_n.\] Inductively \(x_n\geq(n+1)/2>1\), the tails \(T_n\) decrease to zero, and \(a_n=O(\log n)\). Let \(K\) be a nonnegative integer random variable with \(\Pr(K\geq n)=T_n\), and put \[N_n=(n-K)_+=\sum_{j=1}^n1_{\{K<j\}},\qquad
 \Psi_n=\mathbb EN_n=n-a_n.\] These are nondecreasing integer counts of events and \(\Psi_n\to\infty\). Writing \(Y_n=\min(K,n)\), the tail formulas give \(\mathbb EY_n=a_n\) and \(\mathbb EY_n^2=\sum_{j=1}^n(2j-1)T_j\). Hence, for \(n\geq2\), \[\operatorname{Var}(Y_n)-\operatorname{Var}(Y_{n-1})
 =(2x_n-1)T_n-T_n^2=1-T_n
 =\Psi_n-\Psi_{n-1},\] where \(2x_nT_n-T_n^2=1\). Since \(\operatorname{Var}(Y_1)=1/4=\Psi_1-1/4\), it follows that \[\operatorname{Var}(N_n)=\Psi_n-\tfrac14
 \quad\hbox{for every }n.\] Yet \(K<\infty\) almost surely and, eventually, \(N_n-\Psi_n=a_n-K=O(\log n)\) almost surely. Thus even this variance asymptotic, expectation normalization and monotone counting structure do not by themselves establish a square-root almost-sure lower fluctuation. No arithmetic realization of this model is asserted.\par
\section*{Supporting checks}
The source PDF SHA256 is \[\text{\scriptsize\texttt{99197c41a34cfac95205d7926a489c6e94595651a9155bbba89652e35aff21f9}}.\] The source-manifest SHA256 is \[\text{\scriptsize\texttt{0a6ddc6778c493c596acb1dcf99a88f0023e955c52ed6b758bf6d1fbf4caea8d}}.\] The batch token is `[internal identifier omitted]`. The scientific work here was performed in this existing conversation with bounded native helpers; backend model equivalence was not independently verified. Human responsibility and approval were not filled in by AI.\par
\end{document}
